\documentclass{article}
\usepackage{csquotes}
\usepackage{hyperref}
\usepackage[a4paper]{geometry}
\usepackage{fancyhdr}
\pagestyle{fancy}
\lhead{Das Delayed-Choice-Experiment}
\rhead{März 2026}
\begin{document}
\section{Das Delayed-Choice-Experiment}
Wird ein \hyperref[Mach-Zehnder-Interferometer]{Mach-Zehnder-Interferometer} aufgebaut, so ist ein Interferenzmuster zu erkennen.
 
Auf beiden Wegen können nun Kristalle eingebaut werden, welche das passierende Photon in zwei teilen; ein Photon, welches weiter zum Schirm geht und ein Photon welches auf ein Detekor gerichtet wird. Passiert dies, verschwindet das Interferenzmuster weil die Photonen als Teilchen agieren müssen weil der Weg festgestellt wurde.
 
Wird diese Information des Wegs aber wieder \textquote{ausradiert} bevor das Photon den Schirm trifft, indem in der Mitte zwischen den Kristallen und den Detektoren ein Strahlenteiler eingebaut wird, so dass Photonen beider Wege auf beide Detektoren treffen können, erscheint das Interferenzmuster erneut.
 
Die \textquote{Entscheidung}, ob sich das Photon nun wie ein Teilchen oder wie eine Welle verhält konnte sogesehen aber erst getroffen werden, nachdem das Photon bereits emittiert wurde. Dies wird auch, weil die Information \textquote{wegradiert} wird, ein \emph{Quantenradierer}-Experiment genannt. 
\end{document}